\SourceLeaf{415}%
use the following expression from \([6]\)
\[\BareRoot{\left\{rr-\left(\frac{g\mp F}{\lambda}\right)^2\right\}}\]
\par
Here too, then, one root gives us two solutions, namely the symmetric positions of points \((1),(2),(3)\) lying on opposite sides of \(\mathrm{(II)}\) at equal distances from it. For \(x=+1\) the positive direction in \(\mathrm{(I)}\) has the same sense as in \(\mathrm{(II)}\), and for \(x=-1\) it has the opposite sense. The sole exception is the case in which, regardless of sign, \(\lambda r=g-F\) (for \(x=+1\)) or \(=g+F\) (for \(x=-1\)) holds. The two symmetric positions of \(\mathrm{(I)}\) then coincide with one another and with line \(\mathrm{(II)}\) itself.

The roots of equation \([7]\) to be excluded clearly include the imaginary roots and the real roots outside the limits \(-1\) and \(+1\), and also the root \(+1\) or \(-1\) itself if \(\lambda r\) is less than \(g-F\) or \(g+F\), respectively, regardless of sign. It can also be proved that at least one of the three roots always belongs to the class of roots that must be excluded, so there can never be more than four distinct solutions in real coordinates. Strictly speaking, \textit{one quite singular} case is an exception to the former statement, in that no root is excluded. This singular case is the one already mentioned in which, for \(x=\pm1\), the ordinate becomes \(=0\) and in which, as is easily proved, the root in question counts twice. That is, the expression to the left of the equality sign in equation \([7]\) contains the factor \((x\mp1)^2\). The equation therefore has only two distinct roots, and the second may indeed be admissible too. This exception thus leaves the conclusion itself intact.

Finally, some of the solutions in real numbers may still be physically inadmissible. Of the circle drawn with center \((1)\) through \((4)\) and \((5)\), only the appropriate arc is the locus of point \((1)\). It is the arc on the positive side of \((4),(5)\) when \(\vartheta\) is less than \(180^0\), and the arc on the negative side when \(\vartheta\) is a reflex angle. The same holds for the other two circles. Our solution so far does not account for this physical condition. Of the various solutions found in real numbers, therefore, only those are admissible in which the values obtained for the abscissas of the three points \((1),(2),(3)\)
